% Source: https://artofproblemsolving.com/community/c3h1932145_geometry_help
% Stable identity: aops-topic-1932145-q01
Take the side of the square to be $4$, so $BP=CQ=1$ and $PC=QD=3$.

\begin{asy}
usepackage("amsmath");
size(5cm);
defaultpen(fontsize(10pt));
pair D=(0,0), C=(4,0), B=(4,4), A=(0,4), P=(4,3), Q=(3,0), R=(48/25,36/25);
fill(D--Q--R--cycle,lightgray);
draw(A--B--C--D--cycle,linewidth(1));
draw(A--Q,linewidth(0.9)); draw(D--P,linewidth(0.9));
pair u=unit(D-R), v=unit(Q-R); real s=0.28;
draw(R+s*u--R+s*(u+v)--R+s*v);
dot(P);dot(Q);dot(R);
label("$A$",A,NW); label("$B$",B,NE); label("$C$",C,SE); label("$D$",D,SW);
label("$P$",P,E); label("$Q$",Q,NE); label("$R$",R,N);
label("$1$",(4,3.5),E); label("$3$",(4,1.5),E);
label("$1$",(3.5,0),S); label("$3$",(1.5,0),S);
label("\mbox{}",(min(currentpicture,user=true)+max(currentpicture,user=true))/2,black); // keeps the text after the figure black when LaTeX typesets the labels inline
\end{asy}

\textbf{Perpendicularity.} Triangles $ADQ$ and $DCP$ have $AD=DC=4$, $DQ=CP=3$ and right angles at $D$ and $C$, so they are congruent (SAS) and $\angle DAQ=\angle CDP$. In triangle $ADR$,
\[
\angle DAR+\angle ADR=\angle CDP+\left(90^\circ-\angle CDP\right)=90^\circ,
\]
so $\angle ARD=90^\circ$, i.e.\ $AQ\perp DP$.

\textbf{Area.} Triangle $ADQ$ has legs $4$ and $3$, hypotenuse $AQ=5$ and area $6$. Triangle $DRQ$ has a right angle at $R$ and shares the angle at $Q$ with triangle $ADQ$, so it is similar to $ADQ$ with ratio $\dfrac{DQ}{AQ}=\dfrac{3}{5}$. Hence
\[
[DQR]=6\cdot\left(\dfrac{3}{5}\right)^2=\dfrac{54}{25},\qquad\dfrac{[DQR]}{[ABCD]}=\dfrac{54}{25\cdot16}=\boxed{\dfrac{27}{200}} .
\]
In particular $[DQR]<[ABCD]$.
